\section{Introducción}

El caso aborda la utilización de un dataset de enfermedad cardíaca como base académica y de prototipo para un proyecto de ciencia de datos en salud pública. El documento organiza las decisiones de D1 a D6 y los supuestos comunes que permiten al resto del equipo desarrollar los análisis de finanzas, planificación, riesgos e integración final.

La estructura se apoya en los contenidos docentes de Gestión de Proyectos de Datos, la pauta de la evaluación, el Plan Maestro del equipo y el caso Heart Disease. Para mantener la trazabilidad, se distingue entre contenido respaldado por estos documentos, evidencia externa de contexto, decisiones de diseño y supuestos de planificación. Las decisiones concretas de arquitectura y los valores de dimensionamiento son adaptaciones para el caso; no se presentan como hechos declarados por una institución hospitalaria real.

El alcance corresponde a un piloto institucional controlado y mantenible. No se dimensiona la infraestructura para las 1.025 filas del dataset académico ni se afirma que este represente la futura población productiva del hospital. El objetivo es disponer de una base seria y comparable para evaluar alternativas, sin simular un despliegue clínico definitivo.

\subsection{Estado y trazabilidad documental}

El informe maestro identifica a Lucas Moncada como responsable del bloque, declara la línea base conceptual D1 a D6 como cerrada en su versión 1.0 y sitúa la evaluación en septiembre de 2026. Su propósito es habilitar los análisis de finanzas, planificación, riesgos e integración final; se evalúa un piloto institucional controlado y mantenible durante un horizonte de 36 meses.

\begin{description}[style=nextline]
  \item[EA1 1.1] Aporta sesgos, ciclo de vida del dato, justicia, transparencia, privacidad, responsabilidad y mitigación.
  \item[EA1 1.2] Aporta viabilidad técnica, datos, seguridad, escalabilidad, riesgos, RBAC, Zero Trust, anonimización y explicabilidad.
  \item[EA1 1.3] Aporta costos cloud, licencias, RRHH, servicios, costo-beneficio, ROI, escenarios, riesgos y mitigaciones.
  \item[EA1 1.4] Aporta la integración técnica, financiera, ética y de riesgos, además de la comunicación de supuestos y limitaciones.
  \item[Rúbrica EP1] Exige ética 30\%, finanzas 30\%, riesgos 40\%, extensión de 4 a 6 páginas y referencias APA.
  \item[Plan Maestro] Define responsabilidades D1 a D22, criterios de aceptación, dependencias y handoffs.
  \item[Caso y dataset] Fundamentan la necesidad de comprender factores asociados al riesgo y priorizar prevención con un dataset académico de referencia.
\end{description}
